\chapter{第 10 章分层答案}
\label{app:ch10-solutions}


\begin{enumerate}[leftmargin=*]
  \item 无偏、一致和渐近正态是三种不同量词。先写 $\mathbb E\hat\theta_n-\theta$、$\Pr(|\hat\theta_n-\theta|>\varepsilon)$ 和 $a_n(\hat\theta_n-\theta)$ 的对象，再检查方差/矩条件；最后才谈同一估计量类中的有效性。例子 $T_n=X_1$ 的偏差为 0 但方差不随 $n$ 降低，直接证明无偏不推出一致。
  \item 大小是 $\sup_{\theta\in\Theta_0}\Pr_\theta(\text{reject})$，功效是固定备择下的拒绝概率，$p$ 值是观测统计量在零假设尾部的概率，覆盖率是重复抽样事件，经济效应则带原始单位。报告先给效应/区间，再给检验量和成本含义，不能用 p 值替代经济重要性。
  \item 同方差 OLS 使用 $\operatorname{Var}(\hat\beta\mid X)=\sigma^2(X^TX)^{-1}$；异方差时改用 HC，时间相关时用 HAC，簇相关时按簇聚合。每种稳健标准误只改变方差估计，不能修复遗漏变量、未来信息、反向因果或错误函数形式；答案应写出其独立性/簇数前提。
  \item 写 $\hat\theta-\theta=(\hat\theta-\mathbb E\hat\theta)+b$，平方并取期望，交叉项因中心化为零，得 $\mathrm{MSE}=\mathrm{Var}+b^2$。收缩估计偏差 $-0.2$、方差 $0.16$，所以 MSE $0.20$；无偏估计方差 $0.25$。这是在同一目标、同一分布下比较，不能用另一目标的偏差/方差拼接。
  \item $Q(b)=y^Ty-2b^TX^Ty+b^TX^TXb$，逐项微分得 $-2X^Ty+2X^TXb=0$。满列秩时解为 $(X^TX)^{-1}X^Ty$，把它代回 $X^T(y-Xb)$ 得 0；若秩亏，解不唯一或需伪逆，正交性仍是最小范数解的残差条件而非唯一参数识别。
  \item 令遗漏变量 $z=\pi x+v$ 且 $\operatorname{Cov}(x,v)=0$。短回归协方差为 $\operatorname{Cov}(x,y)=\beta_x\operatorname{Var}(x)+\beta_z\operatorname{Cov}(x,z)$，除以 $\operatorname{Var}(x)$ 得 $\beta_x+\beta_z\pi$；代入 $1,2,0.5$ 得 2，偏差 1。稳健协方差只影响区间，不改变这个极限。

  \item 固定第一项在已知方差独立正态模型下具有约 95\% 覆盖率；选择 20 项最大者改变抽样分布，普通区间覆盖率下降。

  \item $e=y-X\hat\beta$，HC0 协方差为 $(X^TX)^{-1}X^T\operatorname{diag}(e_i^2)X(X^TX)^{-1}$。它允许异方差，但不纠正遗漏变量导致的系数偏差。

  \item 标准误从 $\sigma/5$ 降为 $\sigma/10$，独立样本量增至四倍使标准误减半。
  \item 若十次检验在全局零假设下独立，每次仍用 5\% 水平，则至少一次误拒的概率为 $1-.95^{10}=.401263$。反复查看后停止需要相应的序贯方法，例如按预先规定的方式分配总体显著性水平；另一种做法是在探索后使用独立的后期样本确认。
  \item HC3 只对异方差方差估计稳健。先排查处理后变量/碰撞点、遗漏暴露、测量误差、非线性、反向因果、幸存者偏差和时间泄漏；若作因果声明，还要给处理、潜在结果、识别、安慰剂和敏感性分析。每个排查对应一个可执行负例。
  \item BH 应对完整的 500 项检验族应用，并满足相应的依赖条件。它约束错误发现比例，不能直接判断交易方向是否稳定或净收益是否为正。后期样本中的效应、换手及成本，需要另行比较。
\end{enumerate}

\section*{基础衔接：独立计算}
半宽从 $1.96(2/5)=0.784$ 降至 $1.96(2/10)=0.392$。这假定观测独立且标准差已知；仅增加相关行数不保证减半。

\section*{从模型构造估计量}
\begin{enumerate}[leftmargin=*]
\item $L(p)=p^7(1-p)^3$，内点导数 $\ell'(p)=(7-10p)/(p(1-p))$，在 $0.7$ 左边为正、右边为负，故最大点为 $0.7$。全零时 $L(p)=(1-p)^{10}$ 在 $[0,1]$ 上递减，最大点是边界 0；不能仅靠内部驻点判断。
\item $X^TX\hat\beta=X^Ty$ 给 $3\hat\beta_0+3\hat\beta_1=5$、$3\hat\beta_0+5\hat\beta_1=6$，所以 $\hat\beta=(7/6,1/2)^T$。残差平方和 $S=1/6$，最大化 $-3\log v/2-S/(2v)$ 得 $\hat v=S/3=1/18$。无偏方差为 $S/(3-2)=1/6$，因为估计两个系数后只有一个残差自由度；似然以样本密度最大为目标，并不补偿这项自由度消耗。
\item 精确对数似然增加 $\log f_\theta(y_1)$。平稳正态 AR(1) 中，该密度的均值 $c/(1-\phi)$ 与方差 $v/(1-\phi^2)$ 都依赖参数，因此它不是可随意删去的常数。条件似然选择给定 $y_1$，是不同的估计目标；不能声称两者有限样本最大点总相同。
\end{enumerate}
